\section{Introduction}
\label{sec:intro}

Let $S$ be finite, and let $M=(m_{st})_{s,t\in S}$ be a symmetric Coxeter
matrix. Thus $m_{ss}=1$ and, for $s\ne t$,
$m_{st}\in\{2,3,\ldots\}\cup\{\infty\}$. The associated Artin group is
\[
 A=A_M=\left\langle \sigma_s\ (s\in S)\ \middle|\
 \underbrace{\sigma_s\sigma_t\sigma_s\cdots}_{m_{st}}
 =\underbrace{\sigma_t\sigma_s\sigma_t\cdots}_{m_{st}}
 \quad(m_{st}<\infty)\right\rangle.
\]
For $X\subseteq S$, write $A_X=\langle\sigma_s:s\in X\rangle$, with
$A_\varnothing=\{1\}$. A \emph{parabolic subgroup} is a conjugate of
some $A_X$. The rank here refers to the cardinality of the defining set;
no finiteness assumption on the associated Coxeter group is understood.
The Coxeter group is obtained by imposing $\sigma_s^2=1$ for all $s$;
a type, or a parabolic of that type, is called \emph{spherical} when
its Coxeter group is finite.

\begin{theorem}[Parabolic intersections]
\label{thm:main}
For every Coxeter matrix $M$ on a finite set $S$, subsets $X,Y\subseteq S$, and
$g,h\in A_M$, there are $Z\subseteq S$ and $k\in A_M$ such that
\[
 gA_Xg^{-1}\cap hA_Yh^{-1}=kA_Zk^{-1}.
\]
\end{theorem}

In particular, the result applies to arbitrary conjugates of nonspherical
parabolics and to diagrams with infinite labels. The passage from pairwise intersections to arbitrary intersections and
parabolic closure is established in
\cite[Proposition~1.5 and Corollary~1.6]{MoellerParisVarghese2024}.
Our inside-parabolic description gives the following form, whose proof
in Section~\ref{sec:consequences} also records a finite witness bound.

\begin{corollary}[Arbitrary intersections and parabolic closure]
\label{cor:arbitrary-intersections}
For every family $\mathcal F$ of parabolic subgroups of $A$, there is a
subfamily $\mathcal F_0\subseteq\mathcal F$ with $|\mathcal F_0|\le |S|$ such
that
\[
 \bigcap_{P\in\mathcal F}P=\bigcap_{P\in\mathcal F_0}P.
\]
Here the intersection of an empty family is $A$. In particular, arbitrary
intersections of parabolic subgroups are parabolic, and every subset
$E\subseteq A$ is contained in a unique smallest parabolic subgroup.
\end{corollary}


Writing $\operatorname{PC}(E)$ for this parabolic closure, the parabolic
subgroups form a complete lattice: their meet is intersection and the
join of a family is the parabolic closure of its union.

The same stabilizer description gives an effective membership test.
Here the conjugating element and the tested element are supplied as
words in the standard generators and their inverses.

\begin{corollary}[Word and parabolic membership problems]
\label{cor:membership}
There is an algorithm which, given a finite Coxeter matrix $M$, a subset
$X\subseteq S$, and words $u,v$ in the standard generators of $A_M$ and their inverses, decides
whether the element represented by $v$ belongs to $uA_Xu^{-1}$.
In particular, the word problem is decidable for every finite-rank
Artin group.
\end{corollary}

The proof in Section~\ref{sec:algorithms} constructs the finite algebra
with exact algebraic coefficients and tests a single transported
projective by Gaussian elimination.

The intersection theorem also supplies the hypothesis of a published
criterion for Artin groups with a visual splitting.  We state the
result with the resulting trivial-intersection witness; its proof is
in Section~\ref{sec:consequences}.

\begin{corollary}[Non-clique Artin-group dynamics]
\label{cor:nonclique-dynamics}
Suppose $M$ is irreducible, meaning that there is no nontrivial partition
$S=S_1\sqcup S_2$ with $m_{st}=2$ for every $s\in S_1$, $t\in S_2$, and
suppose $m_{st}=\infty$ for some pair. Then $A$ is acylindrically hyperbolic.
Every proper parabolic subgroup $P$ has a conjugate meeting it trivially:
there is $g\in A$ such that $P\cap gPg^{-1}=\{1\}$.
\end{corollary}


The acylindrical-hyperbolicity conclusion in this entire non-clique scope
already appears in \cite[Theorem~1.2]{KatoOguni2025}. Corollary~\ref{cor:nonclique-dynamics}
records the deduction through the parabolic-intersection criterion, not a
new acylindrical-hyperbolicity class.

Theorem~\ref{thm:main} also removes the intersection hypothesis from
the fixed-point and locally compact group results of
\cite{MoellerParisVarghese2024}; their precise consequences are stated
in Section~\ref{sec:fixed-point-consequences}.

\subsection{Context and method}

Standard parabolic subgroups retain the Artin presentation of their
submatrix: this unrestricted embedding theorem is due to van der Lek
\cite{VanDerLek1983}, with subsequent treatments in
\cite{Paris1997,CharneyParis2014}.  Charney and Paris also prove convexity of standard parabolics in the
standard word metric \cite[Theorem~1.2]{CharneyParis2014}.
Blufstein and Paris prove the general nesting theorem: an ambient
parabolic subgroup contained in a parabolic subgroup $Q$ is parabolic
as a subgroup of $Q$ \cite[Theorem~1.1]{BlufsteinParis2023}.
The intersection problem asks whether the intersection of two arbitrary
ambient parabolics is itself parabolic.  It is the Parabolic Intersection Conjecture; precise
formulations appear in \cite[Conjecture~1]{Godelle2023} and
\cite[Section~3, Conjecture~4]{Cumplido2026Survey}.
Theorem~\ref{thm:main} resolves it affirmatively.

Several approaches established this closure in substantial classes.
For right-angled Artin groups, Duncan, Kazachkov and Remeslennikov proved
even the arbitrary-family intersection statement
\cite[Proposition~2.8]{DuncanKazachkovRemeslennikov2007}.
In spherical type, Cumplido, Gebhardt, Gonz\'alez-Meneses and Wiest
proved pairwise closure using Garside theory and parabolic closures
\cite[Theorem~9.5]{CumplidoEtAl2019}.  Their argument builds on the
finite-type positive-monoid structure developed by Brieskorn and Saito
and by Deligne \cite{BrieskornSaito1972,Deligne1972}.
Such positive fractions are available inside spherical subgroups;
they are not available as a general premise for the ambient groups
considered here.

Geometric methods provide a different route.  Cumplido, Martin and
Vaskou obtained arbitrary intersections in large type through fixed
paths in the Artin complex and stabilizer induction
\cite[Theorem~A and Theorem~11]{CumplidoMartinVaskou2023}.
Blufstein extended this approach to $(2,2)$-free two-dimensional Artin
groups \cite[Theorem~1.3]{Blufstein2022}.
Haettel proved the result for affine types $\widetilde A$ and
$\widetilde C$ using geodesic bicombings
\cite[Corollary~8.3]{Haettel2024}; Cumplido, Gavazzi and Paris later
gave a short alternative proof in type $\widetilde A$
\cite[Theorem~1]{CumplidoGavazziParis2024}.

The distinction between the ambient group and the two intersected
subgroups is especially important in FC type.  Morris-Wright proved
the result for two spherical parabolics in an FC-type group, using its
CAT$(0)$ Deligne complex \cite[Theorem~3.1]{MorrisWright2021}.
M\"oller, Paris and Varghese showed that one spherical parabolic suffices
\cite[Corollary~1.2]{MoellerParisVarghese2024}.
Antol\'in and Foniqi proved full pairwise closure for even FC-type groups
by retractions and Bass--Serre theory
\cite[Theorem~1.1]{AntolinFoniqi2022}.
The full FC-type question also appears explicitly in
\cite[Question~3.6]{CharneyMartinMorrisWright2026}.
These results distinguish the closure problem for arbitrary conjugates
from the embedding and nesting properties of standard subgroups.

The word problem has a parallel history. Brieskorn and Saito and
Deligne solved it in spherical type through positive normal forms
and the finite-type monoid structure
\cite[Section~6, especially Theorem~6.6]{BrieskornSaito1972},
\cite[Section~4.20]{Deligne1972}.
For nonspherical examples, Holt and Rees construct shortlex normal
forms in large type \cite[Theorem~3.2]{HoltRees2012}.
Blasco-Garc\'ia, Cumplido, Holt, Morris-Wright and Rees give a
quadratic-time reduction to geodesic words when the Coxeter diagram
has no $A_3$ or $B_3$ subdiagram
\cite[Theorem~1.1]{BlascoGarciaEtAl2026}.
Their introduction explicitly poses the general decidability question;
Corollary~\ref{cor:membership} gives a uniform decision procedure with
the finite Coxeter matrix as part of the input. No complexity bound
is asserted here.

Our proof also reduces intersections to stabilizers and then to smaller
parabolic types.  The space on which we measure a stabilizer is supplied
by a categorical action and a relative harmonic height.  The algebraic
construction belongs to the zigzag and braid-action framework of
Huerfano--Khovanov and Khovanov--Seidel
\cite{HuerfanoKhovanov2001,KhovanovSeidel2002}, with fusion-labelled
Artin actions as in Heng--Licata \cite{HengLicata2024}.
The positive-factor tests are related to Brav--Thomas
\cite{BravThomas2011}.  We construct the required bimodule action and prove the layer and local
rigidity statements, using finite spherical tensor categories and the
standard long exact cohomology sequence as background inputs.  In particular the proof does not use
asphericity of a Salvetti complex.

Heng and Licata give a Gaussian-elimination algorithm for the word
problem in the image of their spherical-twist action
\cite[Proposition~7.11]{HengLicata2024}, and formulate faithfulness
of their original Artin-group action as Conjecture~7.10.
Our relative construction instead detects standard-parabolic
membership by a single framing projective. At the level of decidability,
this membership problem also follows from the word problem and
Charney--Paris convexity: for an input word of length $n$, test all
words of length at most $n$ in the specified parabolic generators and
their inverses. A supplied conjugate is handled by first conjugating
the input word.

\subsection{From relative families to stabilizers}

Fix $X\subseteq S$. Adjoin a vertex $f$ which commutes with the colors in
$X$ and has braid label $3$ with the colors in $S\setminus X$.
Section~\ref{sec:category} constructs a finite-dimensional graded
algebra $Z$, an action $x\mapsto B_x$ of $A$ on the homotopy category
of bounded complexes of graded projective $Z$-modules, and a distinguished finite family of
projectives $(P_{f,L})_L$.  The index $L$ ranges over the simple labels
of a finite semisimple tensor category.  Write
$E_L(x)=B_xP_{f,L}$ for the transported family.
For $Y\subseteq S$ put
\[
 H(Y,x)=\{a\in A_Y:E_L(ax)\simeq E_L(x)\text{ for every }L\},
\]
where equivalence means graded chain homotopy equivalence. The two
claims that drive the proof are
\[
 H(S,1)=A_X,
 \qquad
 H(Y,x)\text{ is conjugate inside }A_Y\text{ to a standard parabolic}.
\]
The first claim translates membership in $A_X$ into a condition on the
initial family.  Transport by $x$ then turns the intersection into
\[
 A_X\cap x^{-1}A_Yx=x^{-1}H(Y,x)x.
\]
The second claim therefore supplies the required parabolic type and
conjugator.  The difficulty is to prove that claim when $Y$ is
nonspherical, without using positive fractions in $A_Y$.

We measure this family by counting its projective generators, with
positive weights, separately at each color and each integer level.
The level is the sum of the generator's cochain and internal degrees.
On a spherical color set these counts have a harmonic interpolation:
the prescribed boundary values extend uniquely to a vector annihilated
there by the Coxeter Cartan matrix.  Sections~\ref{sec:layers}--\ref{sec:null}
relate these numerical extensions to actual complexes and cosets.
The key obstruction rules out a nonzero isolated layer whose count
vector is a positive null vector.  The relative frame requires this
obstruction at \emph{every integer level}; connected support supplies
the necessary boundary data.

The resulting height is defined on cosets $A_Tz$, where $T\subseteq Y$
is spherical and $z\in A_Yx$.  Two vertices are joined when the types
and cosets are nested.  Section~\ref{sec:heights} proves that a finite
path between minimum-height vertices can be lowered until it stays
entirely among minima.  Comparable vertices can have equal height,
so this argument must also lower whole tied segments.  Its conclusion
is connectedness of the minimum set.  In spherical type,
Section~\ref{sec:spherical} supplies a different reduction when some
acted test has nonzero cohomology: a sequence of positive operations,
each using every color, cannot have bounded height.
A finite path bound for the algebra turns prolonged vanishing of
tested cohomology into an empty cochain placement, contradicting the
indecomposability of a transported projective.

These two reductions have complementary roles. Harmonic minima either
place a stabilizer in a conjugate of smaller type or isolate a trivially
acting direct factor. If the operated set itself is spherical, however,
a minimum can have that whole set as its type, so harmonic containment
alone need not decrease rank. The positive-height argument gives a
proper-type reduction unless every test is acyclic and the entire group
already fixes the family. Rank induction then closes the argument
without assuming an intersection theorem at any intermediate step.

\subsection{Organization and conventions}

Sections~\ref{sec:category} and~\ref{sec:layers} construct the categorical
family and its layer calculus. Section~\ref{sec:caps} proves cap geometry
and the spherical two-point lemma. Section~\ref{sec:null} proves the
relative null-layer obstruction. Section~\ref{sec:heights} builds the
harmonic filtration and its stabilizer containment;
Section~\ref{sec:spherical} supplies the spherical positive-height
containment. Section~\ref{sec:stabilizers} proves the stabilizer theorem
and Theorem~\ref{thm:main}, then derives the further consequences.

All complexes are bounded and have cochain differential of degree $1$.
Internal grading is separate from cochain grading; the shift $[1]$ lowers
cochain placements by one. A \emph{generator level} is the sum of its
cochain and internal placements. Unless explicitly qualified as ordinary
cochain degree, test heights use the sum of the two degrees. A coset
$A_Tz$ is a right coset. Every order on cosets below requires both
inclusion of types and actual inclusion of cosets. Connectivity of such
a poset always means finite-path connectivity of its undirected
comparability graph.
